% !PRINTED 584
\weberpage{584}

%<<CONT
which we denote, as a function of $u = v/2K$, by $\boldsymbol{s}(u)$. This function then has the double periodicity expressed by
\weq{4}{\begin{aligned}
\boldsymbol{s}(u+1) &= -\boldsymbol{s}(u), \\
\boldsymbol{s}(u+\omega) &= \boldsymbol{s}(u)
\end{aligned}}
All periods of $\boldsymbol{s}(u)$ are contained in the form $2m + n\omega$, in which $m, n$ are integral rational numbers, and they therefore pass over into integers on multiplication by $a$.

Two numbers $u, u'$ that differ only by a period are called congruent with respect to the modulus $2, \omega$:
\[
u' \equiv u \ (\mathrm{mod}\ 2, \omega).
\]
\gesperrt{Since the function $\mathrm{sn}\,v$ takes a value only twice in the period parallelogram, we have}
\[
\boldsymbol{s}(u) = \boldsymbol{s}(u'),
\]
\gesperrt{only when either}
\weq{5}{\begin{aligned}
u &\equiv u' \\
&\text{or} \\
u &\equiv 1 - u'
\end{aligned}\quad (\mathrm{mod}\ 2, \omega)}
\gesperrt{holds, and $\boldsymbol{s}(u)$ vanishes only when}
\[
u \equiv 0, 1 \ (\mathrm{mod}\ 2, \omega)
\]
\gesperrt{holds.}

Since every fractional number in $\Omega$ can be transformed into an integer by multiplication with an integral rational number, the following theorem results from \S\,156, (6):

1. \gesperrt{If $\eta$ is any fractional number of the field $\Omega$, then $\boldsymbol{s}(\eta)$ is an algebraic number, and if in particular $\eta$ can be represented in such a way that its denominator is relatively prime to 2, then $\boldsymbol{s}(\eta)$ is an integral algebraic number.}

Now let
\weq{6}{\mu = y + x\sqrt{\varDelta}}
be an integer in $\Omega$, in which $x, y$ are integral rational numbers\footnote{In the case of an even $\varDelta$, $\mu$ can accordingly be any integer from $\Omega$. In the case of an odd $\varDelta$, $\mu$ is a number of the order [2].}, and
\[
\mu' = y - x\sqrt{\varDelta},
\]
the number conjugate to $\mu$. Further,
\weq{7}{m = \mu\mu' = y^2 - \varDelta x^2}
is the norm of $\mu$, which we assume to be \gesperrt{odd}.
